\subsection{Derivative of Logarithm Function}
\begin{definition}
    [Logarithm Function] 
    The logarithm of a number is the exponent by which another fixed value, the base, must be raised to produce that number. 
    It can also be recognised as the inverse of an exponential function. 
    The function is always represented in this form:
    \[
        y = \log_{b}x
    \]
    which b is the base, x is the number, y is the exponent. 
\end{definition}

\begin{proposition}
    Given function $y = \log_b x$, the derivative of this function $\frac{d}{dx} \left(\log_b x\right) = \frac{1}{x\ln b}$
    \begin{proof}
        First, we rewrite the function:
        \begin{align}
            y &= \log_{b} x \nonumber\\
            b^y &= x \label{eq5.1.1}
        \end{align}
        Then, we take the derivative of both sides
        \begin{align*}
            b^y \ln b \frac{dy}{dx} &= 1 \\
            \frac{dy}{dx} &= \frac{1}{b^y (\ln b)} 
        \end{align*}
        Sub \ref{eq5.1.1} into the equation
        \begin{align*}
            \frac{dy}{dx} &= \frac{1}{x\ln b}
        \end{align*}
    \end{proof}
\end{proposition}

\begin{example}
    Find any critical numbers for the function $f(x) = (x + 2)e^{-2x}$.
    Classify any critical points. 
    \begin{align*}
        f'(x) &= e^{-2x} + (x + 2)e^{-2x}(-2) \\
        &= e^{-2x}\left(1 - 2x - 4\right) \\
        &= e^{-2x}\left(-2x - 3\right)
    \end{align*}
\end{example}
